\documentclass[a4paper, 14px, fleqn]{article}

% Packages importés
\usepackage[utf8]{inputenc} % Encodage des entrées
\usepackage[T1]{fontenc} % Encodage de la police
\usepackage[french]{babel} % Support pour la langue
\usepackage{fancyhdr} % En-têtes et pieds de page
\usepackage{graphicx} % Ajouter des images
\usepackage{hyperref} % Ajouter les liens dans la table de matière
\usepackage{colortbl} % Ajouter des couleurs dans un tableau
\usepackage{amsmath,amssymb,polynom} % Ajouter les maths
\usepackage{mathtools} % Ajouter des maths
\usepackage{tcolorbox} % Ajouter les cadres
\usepackage{stmaryrd} % Ajouter les doubles brackets
\usepackage{tikz} % Ajouter les schémas
\usetikzlibrary{matrix,arrows,decorations.pathmorphing}
\usepackage{verbatim} % Ajouter les matrices
\usepackage{nicematrix} % Ajouter les matrices
\usepackage[overload]{empheq}

% Utiliser de bonnes marges
\setlength{\hoffset}{-18px}
\setlength{\oddsidemargin}{0px} % Marge gauche sur pages impaires
\setlength{\evensidemargin}{9px} % Marge gauche sur pages paires
\setlength{\marginparwidth}{54px} % Largeur de note dans la marge
\setlength{\textwidth}{481px} % Largeur de la zone de texte (17cm)
\setlength{\voffset}{-18px} % Bon pour DOS
\setlength{\marginparsep}{7px} % Séparation de la marge
\setlength{\topmargin}{0px} % Pas de marge en haut
\setlength{\headheight}{13px} % Haut de page
\setlength{\headsep}{10px} % Entre le haut de page et le texte
\setlength{\footskip}{27px} % Bas de page + séparation
\setlength{\textheight}{708px} % Hauteur de la zone de texte (25cm)

% Ajouter les ensembles
\newcommand{\ensR}{\mathbb{R}}
\newcommand{\ensC}{\mathbb{C}}
\newcommand{\ensK}{\mathbb{K}}
\newcommand{\ensN}{\mathbb{N}}

% Theoremes
\usepackage{tikz, tkz-tab}
\usepackage[framemethod=tikz]{mdframed}
\usetikzlibrary{shapes.misc}

% Les arguments des theoremes importants
\makeatletter
\def\mdf@@title{}
\define@key{mdf}{title}{\def\mdf@@title{#1}}
\def\mdf@@color{}
\define@key{mdf}{color}{\def\mdf@@color{#1}}

% Le style des theoremes
\mdfdefinestyle{mytheo}{
    linewidth=1pt,
    innertopmargin=1.5\baselineskip,
    innerbottommargin=0.7\baselineskip,
    roundcorner=5pt,
    backgroundcolor=\mdf@@color!05,
    linecolor=\mdf@@color,
    singleextra={
        \node[xshift=10pt,thick,draw=\mdf@@color,fill=\mdf@@color!65,rounded 
            corners,text=white,anchor=west] at (P-|O) %
        {\strut{\bfseries\mdf@@title}};
    }
}

% Pour definir le theme comme servant d'environement
\newmdenv[style=mytheo]{theor}

% L'environement mais cette fois il peut prendre les parametres
\newenvironment{theorem}[2][]
  {\begin{theor}[color=#1, title=#2]}
  {\end{theor}}

\makeatother


% Exemples
\newcommand{\exobox}[1]
{
    \begin{tcolorbox}
    {#1}
    \end{tcolorbox}
}

% Mettre en place les en-têtes et les pieds de pages
\pagestyle{fancy}
\fancyhf{}
\lhead{BackToBasics}
\rhead{Mathématiques S3}
\lfoot{\includegraphics[scale=0.15]{img/b2b.png}}
\cfoot{\thepage}
\rfoot{\includegraphics[scale=0.01]{img/epita.png}}

\begin{document}

    % Première page du document
    \vspace*{\stretch{1}}
\begin{center}

    \vspace{15px}
    \Huge{\textbf{Mathématiques}}\\
	\huge{\textbf{Séries numériques}}\\
	\vspace{10px}
	\Large{\textbf{BackToBasics}}\\
	\textit{6 mars 2022}

    \vspace{25px}
	\includegraphics{img/b2b.png}
	
\end{center}
\vspace*{\stretch{1}}
    \clearpage
    
    % Page de table de matières
    \tableofcontents
    \clearpage
    
    % Page de cours
    \section{Cours}
    \noindent Dans ce cours, nous allons revoir tout ce qui concerne les séries numériques. Si vous avez des questions sur n'importe quel point du cours, n'hésitez pas à poser votre question dans le salon dédié ou par messages privés sur Discord à \textbf{Swarwerth\#2943}.

\subsection{Définitions générales}
\begin{theorem}[red]{Définition 1.1 : Série}
    Soit $(u_n)$ une suite réelle. La \textbf{série} de terme générale $u_n$ notée $\sum u_n$ est la suite des \textit{sommes partielles} $(S_n)$ où pour tout $n \in \ensN$, $\displaystyle S_n = \sum_{k = 0}^{n} u_k$. On dit que $\sum u_n$ \textit{converge} si $(S_n)$ converge. On dit qu'elle \textit{diverge} sinon.
\end{theorem}

\begin{theorem}[black]{Exemple de la série géométrique}
    Soit $q \in \ensR$. Alors, $\sum q^n$ converge ssi $|q| < 1$. \\
    $\sum q^n$ est appelée \textit{série géométrique}.
\end{theorem}

\begin{theorem}[blue]{Propriété 1.2 : Stabilité de la convergence d'une série}
    Soient $\sum u_n$, $\sum v_n$ deux séries numériques et $\lambda \in \ensR$. Alors

    \begin{itemize}
        \item ($\sum u_n$ converge et $\sum v_n$ converge) $\Longrightarrow$ $\sum(u_n + v_n)$ converge.
        \item $\sum u_n$ converge $\Longrightarrow$ $\sum \lambda u_n$ converge.
        \item ($\sum u_n$ converge et $\sum v_n$ diverge) $\Longrightarrow$ $\sum(u_n + v_n)$ diverge.
    \end{itemize}
\end{theorem}

\begin{theorem}[red]{Définition 1.3 : Somme et reste d'une série convergente}
    Soit $\sum u_n$ une série numérique convergente. On appelle \textbf{somme} de la série le nombre réel $S$ et le \textbf{reste} de la série, la suite $(R_n)$ définie pour tout $n \in \ensN$ par

    $$
        S = \sum_{n = 0}^{+ \infty} u_n = \lim_{n \to + \infty} S_n
    $$

    $$
        R_n = \sum_{k = n + 1}^{+ \infty} u_k
    $$
\end{theorem}

\begin{theorem}[blue]{Propriété 1.4 : Télescopage}
    Soit $(u_n)$ une suite réelle. Alors,

    $$
        (u_n) \text{ converge } \Longleftrightarrow \sum (u_{n + 1} - u_n) \text{ converge }
    $$
\end{theorem}

\begin{theorem}[blue]{Propriété 1.5 : Condition nécessaire de convergence}
    Soit $(u_n)$ une suite relle. Alors

    $$
        \sum u_n converge \Longrightarrow \lim_{n \to + \infty} u_n = 0 
    $$
\end{theorem}
\clearpage

\subsection{Séries à termes positifs}
\begin{theorem}[red]{Définition 2.1 : Série à termes positifs}
    On dit qu'une série numérique $\sum u_n$ est \textbf{à termes positifs} si pour tout $n \in \ensN, u_n \geq 0$.
\end{theorem}

\noindent \textbf{Remarque :} Si la série $\sum u_n$ à termes positifs, cela implique que sa suite de ses sommes partielles $(S_n)$ est croissante.

\begin{theorem}[blue]{Propriété 2.2 : Condition de convergence}
    Soient $\sum u_n$ une série à termes positifs et $(S_n)$ la suite de ses sommes partielles. Alors,

    $$
        \sum u_n \text{ converge } \Longleftrightarrow (S_n) \text{ est majorée}
    $$
\end{theorem}

\begin{theorem}[blue]{Propriété 2.3 : Théorème de comparaison}
    Soient $(u_n)$ et $(v_n)$ deux suites réelles telles que pour tout $n \in \ensN$, $0 \leq u_n \leq v_n$. Alors

    \begin{itemize}
        \item $\sum v_n$ converge $\Longrightarrow$ $\sum u_n$ converge.
        \item $\sum u_n$ diverge $\Longrightarrow$ $\sum v_n$ diverge.
    \end{itemize}
\end{theorem}

\begin{theorem}[red]{Définition 2.4 : Série de Riemann}
    On appelle \textbf{série de Riemann} toute série de la forme $\displaystyle \sum \frac{1}{n^\alpha}$ avec $\alpha \in \ensR$.
\end{theorem}

\begin{theorem}[orange]{Théorème 2.5 : Convergence d'une série de Riemann}
    Soit $\alpha \in \ensR$. Alors $\displaystyle \sum \frac{1}{n^\alpha}$ converge ssi $\alpha > 1$.
\end{theorem}

\begin{theorem}[blue]{Proposition 2.6 : Comparaisons de Landau}
    Soient $(u_n)$ et $(v_n)$ deux suite réelles positives.

    \begin{itemize}
        \item Si $u_n = O(v_n)$, alors $\sum v_n$ converge $\Longrightarrow$ $\sum u_n$ converge.
        \item Si $u_n = o(v_n)$, alors $\sum v_n$ converge $\Longrightarrow$ $\sum u_n$ converge.
        \item Si $u_n \sim v_n$, alors $\sum u_n$ et $\sum v_n$ sont de même nature.
    \end{itemize}
\end{theorem}

\begin{theorem}[blue]{Proposition 2.7 : Règle de Riemann}
    Soit $(u_n$ une suite positive. \\
    S'il existe $\alpha > 1$ tel que $\lim_{n \to + \infty} n^\alpha u_n = 0$, alors $\sum u_n$ converge.
\end{theorem}

\clearpage

\begin{theorem}[orange]{Théorème 2.8 : Règle de d'Alembert}
    Soit $(u_n)$ une suite réelle strictement positive telle que

    $$
        \lim_{n \to + \infty} \frac{u_{n + 1}}{u_n} = l \hspace{15px} \text{avec } l \in \ensR_+ \cup \{ + \infty \}
    $$

    \noindent Alors, on a

    \begin{itemize}
        \item $l < 1 \Longrightarrow \sum u_n$ converge
        \item $l > 1 \Longrightarrow \sum u_n$ diverge
    \end{itemize}
\end{theorem}

\begin{theorem}[orange]{Théorème 2.9 : Règle de Cauchy}
    Soit $(u_n)$ une suite réelle strictement positive telle que

    $$
        \lim_{n \to + \infty} \sqrt[n]{u_n} = l \hspace{15px} \text{avec } l \in \ensR_+ \cup \{ + \infty \}
    $$

    \noindent Alors, on a

    \begin{itemize}
        \item $l < 1 \Longrightarrow \sum u_n$ converge
        \item $l > 1 \Longrightarrow \sum u_n$ diverge
    \end{itemize}
\end{theorem}
\clearpage

\subsection{Séries à termes quelconques}
\begin{theorem}[red]{Définition 3.1 : Série alternée}
    Soit $(u_n)$ une suite réelle. On dit que $\sum u_n$ est \textbf{alternée} s'il existe une suite réelle $(a_n)$ positive telle que pour tout $n \in \ensN$, $u_n = (-1)^n a_n$.
\end{theorem}

\begin{theorem}[orange]{Théorème 3.2 : Critère spécial des séries alternées}
    Soit $(u_n)$ une suite réelle alternée. Si $(|u_n|)$ est décroissante et converge vers 0, alors

    \begin{itemize}
        \item $\sum u_n$ converge.
        \item $\forall n \in \ensN, |R_n| \leq |u_{n + 1}|$ \hspace{5px} avec $(R_n)$ la suite des restes associée à $\sum u_n$.
    \end{itemize}
\end{theorem}

\begin{theorem}[black]{Exemple du critère spécial de convergence}
    Soit $\alpha \in \ensR$. Alors, $\displaystyle \sum \frac{(-1)^n}{n^\alpha}$ converge ssi $\alpha > 0$.
\end{theorem}

\begin{theorem}[blue]{Définition 3.3 : Convergence absolue}
    On dit qu'une série numérique $\sum u_n$ \textbf{converge absolument} si la série $\sum | u_n |$ converge.
\end{theorem}

\begin{theorem}[orange]{Théorème 3.4 : Théorème de convergence absolue}
    Soit $\sum u_n$ une suite numérique convergeant absolument. Alors, $\sum u_n$ converge.
\end{theorem}
\clearpage
    \clearpage
    
    % Page d'exercices
    \section{Exercices}
    \noindent Nous avons essayé de mettre l'ensemble des connaissance à savoir pour les Midterms au niveau du chapitre sur les séries numériques. Évidemment, vous avez le droit de relire le cours ou de poser des questions. Ces exercices permettront de pointer les lacunes que vous pourriez avoir.

\subsection{Exercice 1 : Convergence de séries}
\noindent Étudier la convergence des séries $\sum u_n$ de terme général suivant :

\begin{enumerate}
    \item $\displaystyle u_n = \frac{1}{\sqrt{n}} ln(1 + \frac{1}{\sqrt{n}})$.
    \item $\displaystyle u_n = (1 + \frac{1}{n})^n - e$.
    \item $\displaystyle u_n = (\frac{n}{n + 1})^{n^2}$.
    \item $\displaystyle u_n = \frac{(n!)^2}{(2n)!} a^n$ avec $a \in \ensR_+^*$.
    \item $\displaystyle u_n = \frac{(-1)^n ln(n)}{n}$.
    \item $\displaystyle u_n = \frac{n^42^{-n^2} + (-1)^n}{\sqrt{n}}$.
    \item $\displaystyle u_n = ln(1 + \frac{(-1)^n}{2n + 1})$.
\end{enumerate}
\vspace{10px}

\subsection{Exercice 2 : Décomposition de terme général}
\begin{enumerate}
    \item Démontrer que la série $\displaystyle \sum \frac{(-1)^n}{\sqrt{n}}$ converge.
    \item Démontrer l'égalité suivante

    $$
        \frac{(-1)^n}{\sqrt{n} + (-1)^n} = \frac{(-1)^n }{\sqrt{n}} - \frac{1}{n} + \frac{(-1)^n}{n\sqrt{n}} + o(\frac{1}{n\sqrt{n}})
    $$

    \item Étudier la convergence de la série $\displaystyle \frac{(-1)^n}{\sqrt{n} + (-1)^n}$
\end{enumerate}
\vspace{10px}

\subsection{Exercice 3 : Une inconnue en plus}
\noindent Soient $a \in \ensR$ et la suite $(u_n)$ définie pour tout $n \geq 2$ par $\displaystyle u_n = \frac{(-1)^n}{(n + (-1)^n)^a}$. On cherche à discuter de la nature de $\sum u_n$ en fonction de $a$.

\begin{enumerate}
    \item Montrer que $\displaystyle u_n \sim \frac{(-1)^n}{n^a}$.
    \item Montrer que, si $a \leq 0$, la série $\sum u_n$ diverge.
    \item On suppose dans la suite de l'exercice que $a > 0$. Déterminer $\lambda \in \ensR$ tel que $\displaystyle u_n = \frac{(-1)^n}{n^a} + \frac{\lambda}{n^{a + 1}} + o(\frac{1}{n^{a + 1}})$.
    \item En déduire la nature de $\sum u_n$.
\end{enumerate}
\vspace{10px}
    \clearpage
    
    % Page de correction
    \section{Correction}
    \noindent Voici la correction des exercices du PDF. Nous vous invitons à faire les exercices avant de lire la correction. Pour toute question, merci de vous rediriger vers le cours et les professeurs présents.

\subsection{Exercice 1 : Convergence de séries}
\noindent Étudier la convergence des séries $\sum u_n$ de terme général suivant :

\begin{enumerate}
    \item $\displaystyle u_n = \frac{1}{\sqrt{n}} ln(1 + \frac{1}{\sqrt{n}})$.

    \exobox
    {
        \begin{align*}
            u_n &= \frac{1}{\sqrt{n}} ln(1 + \frac{1}{\sqrt{n}}) \\
            &= \frac{1}{\sqrt{n}} (\frac{1}{\sqrt{n}} + o(\frac{1}{\sqrt{n}})) \\
            &= \frac{1}{n} + o(\frac{1}{n}) \\
            &\sim \frac{1}{n}
        \end{align*}
    
        Or, $\forall n \in \ensN^*, u_n > 0$ et $\frac{1}{n} > 0$. \\
        De plus, $\sum \frac{1}{n}$ diverge d'après la série de Riemann. \\
        Donc, par comparaison, $\sum u_n$ diverge.
    }

    \vspace{10px}
    
    \item $\displaystyle u_n = (1 + \frac{1}{n})^n - e$.

    \exobox
    {
        \begin{align*}
            u_n &= (1 + \frac{1}{n})^n - e \\
            &= e^{nln(1 + \frac{1}{n})} - e \\
            &= e^{n(\frac{1}{n} - \frac{1}{2n^2} + o(\frac{1}{n^2}))} - e \\
            &= e^{1 - \frac{1}{2n} + o(\frac{1}{n})} - e \\
            &= e^1 \times e^{- \frac{1}{2n} + o(\frac{1}{n})} - e \\
            &= e(1 - \frac{1}{2n} + o(\frac{1}{n})) - e \\
            &= e - \frac{e}{2n} + o(\frac{1}{n}) - e \\
            &= \frac{-e}{2n} + o(\frac{1}{n}) \\
            &\sim \frac{-e}{2n}
        \end{align*}

        Or, $\forall n \in \ensN^*, u_n < 0$ et $\frac{1}{n} < 0$. \\
        De plus, $\sum \frac{1}{n}$ diverge d'après la série de Riemann. On peut alors dire que la série $\sum \frac{-e}{2n}$ diverge. \\
        Donc, par comparaison, $\sum u_n$ diverge.
    }

    \clearpage

    \item $\displaystyle u_n = (\frac{n}{n + 1})^{n^2}$.

    \exobox
    {
        \begin{align*}
            \sqrt[n]{u_n} &= (\frac{n}{n + 1})^n \\
            &= (\frac{n + 1}{n})^{-n} \\
            &= e^{-nln(\frac{n + 1}{n})} \\
            &= e^{-nln(1 + \frac{1}{n})} \\
            &= e^{-n(\frac{1}{n} + o(\frac{1}{n}))} \\
            &= e^{-1 + o(1)}
        \end{align*}

        Donc on a

        $$
            lim_{n \to + \infty} \sqrt[n]{u_n} = e^{-1} = \frac{1}{e}
        $$

        Or, $\frac{1}{e} < 1$. \\
        De plus, $\forall n \in \ensN, u_n > 0$. \\
        Donc, d'après la règle de Cauchy, \\
        $\sum u_n$ est convergente.
    }

    \clearpage
    
    \item $\displaystyle u_n = \frac{(n!)^2}{(2n)!} a^n$ avec $a \in \ensR_+^*$.

    \exobox
    {
        \begin{align*}
            \frac{u_{n + 1}}{u_n} &= \frac{((n + 1)!)^2}{(2n + 2)!} a^{n + 1} \times \frac{(2n)!}{(n!)^2} \frac{1}{a^n} \\
            &= (\frac{(n + 1)!}{n!})^2 \times \frac{(2n)!}{(2n + 2)!} \times \frac{a^{n + 1}}{a^n} \\
            &= (n + 1)^2 \times \frac{1}{(2n + 2)(2n + 1)} \times a \\
            &= \frac{n^2 + 2n + 1}{4n^2 + 4n + 2n + 2} \times a \\
            &\sim \frac{n^2}{4n^2} \times a \\
            &\sim \frac{a}{4}
        \end{align*}

        Or, $\forall n \in \ensN, u_n > 0$. \\
        Donc, en utilisant la règle de d'Alembert, on a :

        \begin{itemize}
            \item Si $a < 4$, $\sum u_n$ converge.
            \item Si $a > 4$, $\sum u_n$ diverge.
        \end{itemize}

        Si on prend $a = 4$, on a $u_n = \frac{(n!)^2}{(2n)!} \times 4^n$.

        \begin{align*}
            \frac{u_{n + 1}}{u_n} &= \frac{(n + 1)^2}{(2n + 1)(2n + 1)} \times 4 \\
            &= \frac{4n^2 + 8n + 4}{4n^2 + 4n + 2n + 2} \\
        \end{align*}

        Or, $n \in \ensN$, donc $\frac{u_{n + 1}}{u_n} > 1 \Longleftrightarrow u_{n + 1} > u_n$. \\
        Donc, $(u_n)$ est croissante et $\forall n \in \ensN, u_n \geq 0$. \\
        Ainsi, $(u_n)$ ne peut pas tendre vers 0. \\
        On en conclut que $\sum u_n$ diverge.
    }

    \vspace{10px}
    
    \item $\displaystyle u_n = \frac{(-1)^n ln(n)}{n}$.

    \exobox
    {
        La série $(u_n)$ est alternée, \\
        On a alors
    
        $$
            | u_n | = | \frac{(-1)^nln(n)}{n} | = \frac{ln(n)}{n}
        $$

        Par croissance comparée, $(|u_n|)$ est convergente vers 0. \\
        De plus, $(|u_n|)$ est décroissante. \\
        Donc, d'après le critère spécial des séries alternées, $\sum u_n$ converge.
    }

    \clearpage
    
    \item $\displaystyle u_n = \frac{n^42^{-n^2} + (-1)^n}{\sqrt{n}}$.

    \exobox
    {
        $$
            u_n = \frac{n^42^{-n^2} + (-1)^n}{\sqrt{n}} = \frac{n^4 2^{-n^2}}{\sqrt{n}} + \frac{(-1)^n}{\sqrt{n}}
        $$

        Premièrement, nous avons, \\
        $\displaystyle (\frac{(-1)}{\sqrt{n}})$ est une suite alternée. \\
        Donc, on a
        
        $$
            | \frac{(-1)^n}{\sqrt{n}} | = \frac{1}{\sqrt{n}}
        $$

        Or, $\displaystyle (\frac{1}{\sqrt{n}})$ tend vers 0. \\
        De plus, $\displaystyle (\frac{1}{\sqrt{n}})$ est décroissante. \\
        Donc, d'après le critère spécial des séries alternées, $\displaystyle \sum \frac{(-1)^n}{\sqrt{n}}$ converge. \\

        Deuxièmement, nous avons

        \begin{align*}
            \sqrt[n]{\frac{n^4 2^{-n^2}}{\sqrt{n}}} &= \frac{n^{\frac{4}{n}} 2^{-n}}{n^{\frac{1}{2n}}} \\
            &= 2^{-n} \times n^{\frac{4}{n} - \frac{1}{2n}} \\
            &= \frac{1}{2^n} \times n^{\frac{7}{2n}} \\
            &= \frac{e^{\frac{7}{2n}ln(n)}}{e^{nln(2)}}
        \end{align*}

        Or, par croissance comparée, $\displaystyle (\sqrt[n]{\frac{n^4 2^{-n^2}}{\sqrt{n}}})$ converge vers 0. \\
        Donc, d'après la règle de Cauchy, \\
        $\displaystyle \sum \frac{n^4 2^{-n^2}}{\sqrt{n}}$ converge. \\

        Ainsi, par somme de deux séries convergentes, \\
        $\sum u_n$ converge.
    }

    \clearpage
    
    \item $\displaystyle u_n = ln(1 + \frac{(-1)^n}{2n + 1})$.

    \exobox
    {
        \begin{align*}
            u_n &= ln(1 + \frac{(-1)^n}{2n + 1}) \\
            &= \frac{(-1)^n}{2n + 1} - (\frac{(-1)^n}{2n + 1})^2 \times \frac{1}{2} + o(\frac{1}{n^2}) \\
            &= \frac{(-1)^n}{2n + 1} - \frac{(-1)^{2n}}{2(2n + 1)^2} + o(\frac{1}{n^2}) \\
            &= \frac{(-1)^n}{2n + 1} + v_n \hspace{30px} \text{avec } v_n = \frac{1}{2(2n + 1)^2} + o(\frac{1}{n^2})
        \end{align*}

        $$
            v_n = \frac{1}{2(4n^2 + 4n + 1)} + o(\frac{1}{n^2}) \sim \frac{1}{8n^2}
        $$

        Or, $\displaystyle \forall n \in \ensN, \frac{1}{8n^2} > 0$ et $v_n > 0$. \\
        De plus, $\displaystyle \sum \frac{1}{8n^2}$ converge d'après la formule de Riemann. \\
        Donc, $\sum v_n$ converge par comparaison. \\

        $(\frac{(-1)^n}{2n + 1})$ est une suite alternée, \\
        Donc, on a

        $$
            | \frac{(-1)^n}{2n + 1} | = \frac{1}{2n + 1}
        $$

        Or, $\displaystyle (\frac{1}{2n + 1})$ est décroissante et converge vers 0. \\
        Donc, d'après le critère spécial des séries alternées, $\displaystyle \sum \frac{(-1)^n}{2n + 1}$ converge. \\

        Par somme de séries convergentes, \\
        $\sum u_n$ converge.
    }
\end{enumerate}
\clearpage

\subsection{Exercice 2 : Décomposition de terme général}
\begin{enumerate}
    \item Démontrer que la série $\displaystyle \sum \frac{(-1)^n}{\sqrt{n}}$ converge.

    \exobox
    {
        $\displaystyle (\frac{(-1)^n}{\sqrt{n}})$ est une suite alternée. \\
        Donc, on a
    
        $$
            | \frac{(-1)^n}{\sqrt{n}} | = \frac{1}{\sqrt{n}}
        $$

        Or, $\displaystyle (\frac{1}{\sqrt{n}})$ est décroissante et converge vers 0. \\
        D'après le critère spécial des séries alternées, $\displaystyle \sum \frac{(-1)^n}{\sqrt{n}}$ converge.
    }

    \vspace{10px}
    
    \item Démontrer l'égalité suivante

    $$
        \frac{(-1)^n}{\sqrt{n} + (-1)^n} = \frac{(-1)^n}{\sqrt{n}} - \frac{1}{n} + \frac{(-1)^n}{n\sqrt{n}} + o(\frac{1}{n\sqrt{n}})
    $$

    \exobox
    {
        \begin{align*}
            \frac{(-1)^n}{\sqrt{n} + (-1)^n} &= \frac{(-1)^n}{\sqrt{n}(1 + \frac{(-1)^n}{\sqrt{n}})} \\
            &= \frac{(-1)^n}{\sqrt{n}} \times \frac{1}{1 + \frac{(-1)^n}{\sqrt{n}}} \\
            &= X \times \frac{1}{1 + X} \hspace{20px} \text{avec } X = \frac{(-1)^n}{\sqrt{n}} \text{ et } lim_{n \to + \infty} X = 0 \\
            &= X \times (1 - X + X^2 + o(X^2) \\
            &= X - X^2 + X^3 + o(X^3) \\
            &= \frac{(-1)^n}{\sqrt{n}} - \frac{(-1)^{2n}}{(\sqrt{n})^2} + \frac{(-1)^{3n}}{(\sqrt{n})^3} + o(\frac{1}{(\sqrt{n})^3}) \\
            &= \frac{(-1)^n}{\sqrt{n}} - \frac{1}{n} + \frac{(-1)^n}{n\sqrt{n}} + o(\frac{1}{n\sqrt{n}})
        \end{align*}
    }

    \clearpage

    \item Étudier la convergence de la série $\displaystyle \frac{(-1)^n}{\sqrt{n} + (-1)^n}$

    \exobox
    {
        On pose $u_n = \frac{(-1)^n}{n\sqrt{n}} + o(\frac{1}{n\sqrt{n}})$. \\
        $(u_n)$ est une suite alternée, \\
        Donc, on a :

        $$
            | u_n | = \frac{1}{n\sqrt{n}} + o(\frac{1}{n\sqrt{n}}) \sim \frac{1}{n\sqrt{n}}
        $$

        D'après la formule de Riemann, \\
        $\sum \frac{1}{n\sqrt{n}}$ converge. \\
        Donc, par comparaison $\sum |u_n|$ converge. \\
        Ainsi, $\sum u_n$ converge absolument. \\
        On en conclut que $\sum u_n$ converge. \\

        Or, on a $\displaystyle \frac{(-1)^n}{\sqrt{n} + (-1)^n} = \frac{(-1)^n}{\sqrt{n}} - \frac{1}{n} + u_n$. \\

        Par somme de séries convergentes,
        $\displaystyle \sum \frac{(-1)^n}{\sqrt{n}} + u_n$ converge. \\

        Par somme d'une série convergente et d'une série divergente, \\
        $\displaystyle \sum \frac{(-1)^n}{\sqrt{n} + (-1)^n}$ diverge.
    }
\end{enumerate}
\clearpage

\subsection{Exercice 3 : Une inconnue en plus}
\noindent Soient $a \in \ensR$ et la suite $(u_n)$ définie pour tout $n \geq 2$ par $\displaystyle u_n = \frac{(-1)^n}{(n + (-1)^n)^a}$. On cherche à discuter de la nature de $\sum u_n$ en fonction de $a$.

\begin{enumerate}
    \item Montrer que $\displaystyle u_n \sim \frac{(-1)^n}{n^a}$.

    \exobox
    {
        \begin{align*}
            u_n &= \frac{(-1)^n}{(n + (-1)^n)^a} \\
            &= (-1)^n \times \frac{1}{(n(1 + \frac{(-1)^n}{n})^a} \\
            &= \frac{(-1)^n}{n^a} \times (1 + \frac{(-1)^n}{n})^{-a} \\
            &= \frac{(-1)^n}{n^a} (1 + o(1)) \\
            &= \frac{(-1)^n}{n^a} + o(\frac{1}{n^a}) \\
            &\sim \frac{(-1)^n}{n^a}
        \end{align*}
    }

    \vspace{10px}

    \item Montrer que, si $a \leq 0$, la série $\sum u_n$ diverge.

    \exobox
    {
        $$
            |u_n| \sim \frac{1}{n^a}
        $$

        Or, quand $a \leq 0$, $\displaystyle \lim_{n \to + \infty} \frac{1}{n^a} \neq 0$. \\
        Par comparaison de deux suites de terme général positif, \\
        $(|u_n|)$ ne converge pas vers 0. \\
        Donc, $(u_n)$ ne converge pas vers 0. \\
        Ainsi, $\sum u_n$ diverge.
    }

    \vspace{10px}
    
    \item On suppose dans la suite de l'exercice que $a > 0$. Déterminer $\lambda \in \ensR$ tel que $\displaystyle u_n = \frac{(-1)^n}{n^a} + \frac{\lambda}{n^{a + 1}} + o(\frac{1}{n^{a + 1}})$.

    \exobox
    {
        \begin{align*}
            u_n &= \frac{(-1)^n}{(n + (-1)^n)^a} \\
            &= (-1)^n \times \frac{1}{(n(1 + \frac{(-1)^n}{n})^a} \\
            &= \frac{(-1)^n}{n^a} \times (1 + \frac{(-1)^n}{n})^{-a} \\
            &= \frac{(-1)^n}{n^a} (1 - a \times \frac{(-1)^n}{n} + o(\frac{1}{n})) \\
            &= \frac{(-1)^n}{n^a} - a \times \frac{(-1)^{2n}}{n^{a + 1}} + o(\frac{1}{n^{a + 1}}) \\
            &= \frac{(-1)^n}{n^a} + \frac{- a}{n^{a + 1}} + o(\frac{1}{n^{a + 1}})
        \end{align*}

        Donc, on a $\lambda = -a$.
    }

    \clearpage
    
    \item En déduire la nature de $\sum u_n$.

    \exobox
    {
        Premièrement, nous avons $\displaystyle (\frac{(-1)^n}{n^a})$ qui est une suite alternée. \\
        Donc, on a
        
        $$
            | \frac{(-1)^n}{n^a} | = \frac{1}{n^a}
        $$

        Or, $\displaystyle (\frac{1}{n^a})$ est décroissante et tend vers 0 pour tout $a > 0$. \\
        Donc, d'après le critère spécial des séries alternées, $\displaystyle \sum \frac{(-1)^n}{n^a}$ converge. \\

        Deuxièmement, on a

        $$
            \frac{- a}{n^{a + 1}} + o(\frac{1}{n^{a + 1}}) \sim \frac{- a}{n^{a + 1}}
        $$

        Or, $\displaystyle \forall n \in \ensN, \frac{- a}{n^{a + 1}} < 0$. \\
        De plus, d'après la formule de Riemann, $\displaystyle \sum \frac{- a}{n^{a + 1}}$ converge pour tout $a > 0$. \\
        Donc, $\displaystyle \sum \frac{- a}{n^{a + 1}} + o(\frac{1}{n^{a + 1}})$ converge par comparaison. \\

        Par somme de séries convergentes, \\
        $\sum u_n$ converge.
    }
\end{enumerate}
\clearpage
    \clearpage


\end{document}
